\chapter{How to read this blueprint}

This blueprint follows mathematical dependencies rather than the page order of
Dietrich Braess's \emph{Finite Elements: Theory, Fast Solvers, and Applications
in Elasticity Theory}, third English edition (2007). Statements are paraphrased;
the edition, section, and page should be supplied when a node is refined into a
formalization issue.

A node with no \texttt{\textbackslash leanok} annotation is not yet formalized.
The presence of a Lean declaration link does not replace human semantic review.
Each mathematical PR records whether its definitions and statement have been
checked independently. GitHub Projects tracks work flow; this blueprint tracks
mathematical dependencies and formalization status.

\chapter{Abstract variational and Galerkin theory}

The immediate reference is Braess, Chapter II, Section 4. This chapter is the
first milestone and deliberately has no mesh or concrete PDE dependency.

\begin{definition}[Continuous bilinear form]
  \label{def:continuous-bilinear-form}
  \lean{FemLean.Core.Variational.BoundedBy,FemLean.Core.Variational.boundedBy_iff_opNorm_le}
  Let $V$ be a real Hilbert space. A bilinear form $a:V\times V\to\mathbb R$
  is continuous if there is a finite constant $M$ such that
  $|a(u,v)|\le M\|u\|_V\|v\|_V$ for all $u,v\in V$.
  The Lean predicate takes the chosen nonnegative bound as an argument and
  applies to real normed spaces. For a bundled continuous bilinear form,
  its pointwise estimate is equivalent to the operator-norm bound.
\end{definition}

\begin{definition}[Coercive bilinear form]
  \label{def:coercive-bilinear-form}
  \lean{FemLean.Core.Variational.CoerciveWith,FemLean.Core.Variational.isCoercive_iff_exists_coerciveWith}
  A bilinear form $a$ on $V$ is coercive if there is $\alpha>0$ such that
  $\alpha\|v\|_V^2\le a(v,v)$ for every $v\in V$.
  The Lean predicate takes $\alpha$ as an argument. Existential quantification
  over this argument recovers Mathlib's coercivity predicate.
  Symmetry is not assumed unless a later theorem actually needs it.
\end{definition}

\begin{definition}[Weak solution]
  \label{def:weak-solution}
  \uses{def:continuous-bilinear-form}
  Given a continuous linear functional $\ell$ on $V$, a weak solution is an
  element $u\in V$ satisfying $a(u,v)=\ell(v)$ for every $v\in V$.
\end{definition}

\begin{definition}[Galerkin solution]
  \label{def:galerkin-solution}
  \uses{def:continuous-bilinear-form}
  Let $V_h$ be a closed subspace of $V$. A Galerkin solution is $u_h\in V_h$
  satisfying $a(u_h,v_h)=\ell(v_h)$ for every $v_h\in V_h$.
  The representation of the discrete subspace is part of the M1 interface review.
\end{definition}

\begin{lemma}[Galerkin orthogonality]
  \label{lem:galerkin-orthogonality}
  \uses{def:weak-solution,def:galerkin-solution}
  If $u$ is a weak solution and $u_h$ is the corresponding Galerkin solution on
  $V_h$, then $a(u-u_h,v_h)=0$ for every $v_h\in V_h$.
\end{lemma}

\begin{theorem}[Céa's quasi-optimality theorem]
  \label{thm:cea}
  \uses{def:continuous-bilinear-form,def:coercive-bilinear-form,
    def:weak-solution,def:galerkin-solution,lem:galerkin-orthogonality}
  Assume that $a$ is continuous with constant $M$ and coercive with constant
  $\alpha>0$. If $u$ is the weak solution and $u_h\in V_h$ is a Galerkin
  solution, then
  \[
    \|u-u_h\|_V\le \frac{M}{\alpha}
      \inf_{v_h\in V_h}\|u-v_h\|_V.
  \]
  Existence and uniqueness may be obtained through a separate Lax--Milgram
  interface. Reviewers must check that symmetry has not been added accidentally.
\end{theorem}

\begin{corollary}[Exactness on the whole space]
  \label{cor:galerkin-whole-space}
  \uses{thm:cea}
  When $V_h=V$, a Galerkin solution agrees with the weak solution. This is an
  early compiled interface test, not merely an illustrative remark.
\end{corollary}

\chapter{Poisson model problem}

The reference is Braess, Chapter II, Sections 1--4. Work here is blocked until
RFC 0001 selects a Sobolev representation.

\begin{definition}[Poisson weak problem]
  \label{def:poisson-weak}
  \uses{def:weak-solution}
  \notready
  On a bounded domain $\Omega$ with explicitly stated regularity and boundary
  hypotheses, define the weak problem associated with $-\Delta u=f$ and
  homogeneous Dirichlet boundary conditions. The precise choices of $f$, test
  space, and boundary representation belong to RFC 0001 and its issues.
\end{definition}

\begin{proposition}[Poisson Galerkin quasi-optimality]
  \label{prop:poisson-galerkin-cea}
  \uses{def:poisson-weak,thm:cea}
  \notready
  The Poisson bilinear form satisfies the abstract continuity and coercivity
  hypotheses, so every conforming Galerkin approximation inherits Céa's
  quasi-optimality estimate.
\end{proposition}

\chapter{Two-dimensional affine $P_1$ elements}

The references are Braess, Chapter II, Sections 4--8, and RFCs 0001, 0002
and 0004 for the Sobolev, mesh and finite element representations.

\begin{definition}[Triangulation of a fixed domain]
  \label{def:fixed-domain-triangulation}
  \notready
  Fix a bounded polygonal Lipschitz domain $\Omega$ in the Euclidean plane.
  Let $S_h$ be the union of the closed nondegenerate triangles of a finite
  conforming mesh. Require $S_h=\overline{\Omega}$ and
  $\operatorname{int}(S_h)=\Omega$. Distinct triangles meet in a common
  face or have empty intersection. The edges incident to just one triangle
  form the boundary of $\Omega$.
\end{definition}

\begin{definition}[Shape-regular triangular mesh family]
  \label{def:shape-regular-mesh-family}
  \uses{def:fixed-domain-triangulation}
  \notready
  A family of triangulations of the same $\Omega$ has shape bound $\kappa$
  if $h_T/\rho_T\le\kappa$ for every triangle $T$ of every mesh, where
  $h_T$ is its diameter and $\rho_T$ its inradius. For each mesh set
  $h=\max_T h_T$.
\end{definition}

\begin{definition}[Affine $P_1$ finite element]
  \label{def:affine-p1}
  \notready
  On a nondegenerate reference triangle the three vertex values uniquely
  determine an affine polynomial. Transport the nodal basis to each physical
  triangle by affine pullback. The basis satisfies the Kronecker-delta
  property at its vertices.
\end{definition}

\begin{proposition}[$P_1$ conformity]
  \label{prop:p1-conformity}
  \uses{def:affine-p1,def:fixed-domain-triangulation}
  \notready
  Local affine functions reconstructed from the same global nodal data agree
  on each common edge, including when the two local vertex orders are reversed.
  The reconstruction $R_h$ takes arbitrary vertex values into $H^1(\Omega)$.
  Its weak gradient equals the cellwise affine gradient almost everywhere.
\end{proposition}

\begin{proposition}[Homogeneous Dirichlet reconstruction]
  \label{prop:p1-dirichlet-reconstruction}
  \uses{prop:p1-conformity}
  \notready
  Let $D_h^0$ be the coordinate space indexed by the interior vertices.
  Extending its coordinates by zero at the boundary vertices and reconstructing
  defines a linear injection $R_h^0:D_h^0\hookrightarrow H_0^1(\Omega)$.
  Write $V_h^0=\operatorname{range}(R_h^0)$ for its image.
\end{proposition}

\begin{proposition}[Galerkin equations in coordinates]
  \label{prop:galerkin-coordinates}
  \uses{def:galerkin-solution}
  \notready
  Let $(\phi_j)$ be a basis of a finite-dimensional trial and test space and
  write $RU=\sum_j U_j\phi_j$. Define $K_{ij}=a(\phi_j,\phi_i)$ and
  $F_i=\ell(\phi_i)$. Then $a(RU,RW)=W^{\mathsf T}KU$ for all $U,W$, and
  $KU=F$ holds exactly when $RU$ satisfies the Galerkin equation. This
  coordinate identity also applies to nonsymmetric forms.
\end{proposition}

\begin{proposition}[Reduced Poisson system]
  \label{prop:poisson-reduced-system}
  \uses{def:poisson-weak,prop:p1-dirichlet-reconstruction,prop:galerkin-coordinates}
  \notready
  For the unit-coefficient Dirichlet Poisson problem, assemble the element
  integrals and eliminate the boundary coordinates. Then $K_0 U=F_0$ if and
  only if $R_h^0 U$ solves the Galerkin equation on $V_h^0$. The Dirichlet
  form is symmetric and coercive on $H_0^1(\Omega)$. Its restriction along
  the injective $R_h^0$ gives a positive-definite reduced matrix.
\end{proposition}

\begin{proposition}[Boundary-preserving first-order approximation]
  \label{prop:p1-approximation}
  \uses{def:shape-regular-mesh-family,prop:p1-dirichlet-reconstruction}
  \notready
  Fix $\Omega$ and a shape bound $\kappa$. There is a constant
  $C_{\mathrm{app}}(\Omega,\kappa)$ such that, for every mesh with that bound
  and every $u\in H^2(\Omega)\cap H_0^1(\Omega)$, there is a $w_h\in V_h^0$
  with $\|u-w_h\|_{H^1(\Omega)}\le C_{\mathrm{app}} h |u|_{H^2(\Omega)}$.
  For nodal interpolation, evaluate the continuous representative of $u$
  on the closed triangles and justify its zero boundary values.
\end{proposition}

\begin{theorem}[First-order energy error estimate]
  \label{thm:p1-energy-error}
  \uses{prop:poisson-galerkin-cea,prop:p1-approximation}
  \notready
  For the unit-coefficient Dirichlet Poisson problem, fix $\Omega$ and
  $\kappa$. There exists $C=C(\Omega,\kappa)$ such that for every mesh
  with this shape bound and every weak solution
  $u\in H^2(\Omega)\cap H_0^1(\Omega)$, the Galerkin solution on $V_h^0$
  satisfies
  \[
    \|u-u_h\|_{H^1(\Omega)}\le C h |u|_{H^2(\Omega)}.
  \]
\end{theorem}

\begin{corollary}[First-order error for the reduced matrix solution]
  \label{cor:p1-reduced-matrix-error}
  \uses{prop:poisson-reduced-system,thm:p1-energy-error}
  \notready
  With the same hypotheses and the same constant $C$, any solution of
  $K_0 U=F_0$ satisfies
  \[
    \|u-R_h^0 U\|_{H^1(\Omega)}\le C h |u|_{H^2(\Omega)}.
  \]
\end{corollary}
